\documentclass[../../../include/open-logic-section]{subfiles}

\begin{document}

\olfileid{sth}{choice}{woproblem}
\olsection{সুক্রমবিন্যাসের সমস্যা}

সুক্রমবিন্যাস গ্রহণ করি কি না, তার উপর অনেক কিছু
নির্ভর করছে। কিন্তু এই নীতি নিয়ে আলোচনা প্রধানত
সমতুল্য আরেক নীতি—চয়নের স্বতঃসিদ্ধ—নিয়ে হয়েছে।
তাই এবার সেদিকে মন দিই এবং সমতুল্যতাটিও প্রমাণ করি।

\citeyear{Cantor1883}-এ কান্টর সুক্রমবিন্যাসের
স্বতঃসিদ্ধের সমর্থনে একে ``আমার কাছে মৌলিক, ফলে
সমৃদ্ধ, এবং বিশেষত সার্বিক প্রযোজ্যতায় লক্ষণীয় এক
চিন্তার নিয়ম'' বলেছিলেন (\citeauthor{Potter2004}
\citeyear[p.~243]{Potter2004}-এ উদ্ধৃত)। কিন্তু
শেষ পর্যন্ত কান্টর নিজেই বুঝেছিলেন, এই ``স্বতঃসিদ্ধ''
প্রমাণের দাবি রাখে। গণিতসমাজেরও তাই মনে হয়েছিল।

\citeyear{Zermelo1904}-এ জার্মেলো সমস্যাটির
``সমাধান'' করেন। তাঁর সমাধান বুঝতে কয়েকটি সংজ্ঞা দরকার।

\begin{defn}
একটি অপেক্ষক $f$ \emph{চয়ন-অপেক্ষক} তখনই, যখন
প্রত্যেক $x \in \dom{f}$-র জন্য $f(x) \in x$।
$f$-কে \emph{$A$-র চয়ন-অপেক্ষক} বলি তখনই,
যখন $f$ চয়ন-অপেক্ষক এবং
$\dom{f} = A \setminus \{\emptyset\}$।
\end{defn}

স্বজ্ঞায়, $A$-র প্রত্যেক অশূন্য সদস্য-সেট
$x \in A$ থেকে একটি নির্দিষ্ট সদস্য $f(x)$
\emph{বেছে নেয়} $A$-র চয়ন-অপেক্ষক,
সেই $x$ থেকেই।
তাহলে চয়নের স্বতঃসিদ্ধ:

\begin{axiom}[চয়ন]
	প্রত্যেক সেটের একটি চয়ন-অপেক্ষক আছে।
\end{axiom}

জার্মেলো দেখিয়েছিলেন, চয়ন থেকে সুক্রমবিন্যাস
পাওয়া যায়, আবার উল্টোটিও:

\begin{thm}[$\ZF$-এ]\ollabel{thmwochoice}
সুক্রমবিন্যাস ও চয়ন সমতুল্য।
\end{thm}

\begin{proof}
\emph{বাঁ দিক থেকে ডান দিকে।} $A$ সেটের একটি
সেট হোক। সংযোগের স্বতঃসিদ্ধে $\bigcup A$ আছে।
সুক্রমবিন্যাস ধরে $\bigcup A$-র উপর একটি সুক্রম $<$ নিই।
এবার $f(x) = \text{< ক্রমে }x\text{-এর ক্ষুদ্রতম সদস্য}$
নিই। এটি $A$-র চয়ন-অপেক্ষক।

\emph{ডান দিক থেকে বাঁ দিকে।} $A$ স্থির করি।
% BN-SRC-467: the source initial value f(A) is undefined when A is empty.
$A=\emptyset$ হলে সুক্রমবিন্যাস স্পষ্ট। তাই
$A$ অশূন্য ধরি। চয়ন অনুযায়ী
$\Pow{A} \setminus \{\emptyset\}$-র একটি
চয়ন-অপেক্ষক $f$ আছে। ট্রান্সফাইনাইট পুনরাবৃত্তি
দিয়ে সংজ্ঞায়িত করি:
\begin{align*}
	g(0) &= f(A)\\
	g(\alpha) &=
		\begin{cases}
			\text{থামো!{}} &\text{যদি }A = \funimage{g}{\alpha}\\
			f(A \setminus \funimage{g}{\alpha}) & \text{অন্যথায়}\\	
		\end{cases}
\end{align*}
``থামো!'' প্রকৃত সংজ্ঞার সংক্ষিপ্ত রূপ।
% BN-SRC-468: the sentinel is assigned from the first stop onward,
% not retroactively at all earlier stages; injectivity concerns the pre-stop segment.
প্রথমবার $A=\funimage{g}{\alpha}$ হলে সেই
পর্যায় থেকে পরের সব $\delta\geq\alpha$-তে
$g(\delta)=A$
ধরি; আগের মানগুলি অক্ষত থাকে। প্রথমে এতে
শুধু একটি \emph{পদ} সংজ্ঞায়িত হয়েছে বলা যায়
(দেখো \olref[sth][spine][recursion]{transrecursionschema}-এর
পরের মন্তব্য); কিন্তু থামার আগের অংশটি যে
অপেক্ষক, তা দেখাব।

$f$ চয়ন-অপেক্ষক বলে, থামার আগের প্রত্যেক
$\alpha$-র জন্য
$g(\alpha) = f(A \setminus \funimage{g}{\alpha})
\in A \setminus \funimage{g}{\alpha}$;
অতএব $g(\alpha) \notin \funimage{g}{\alpha}$।
এই অংশে $g(\alpha)=g(\beta)$ হলে
$g(\beta)\notin\funimage{g}{\alpha}$,
তাই $\beta\notin\alpha$; একইভাবে
$\alpha\notin\beta$। ক্রমসংখ্যার ত্রিবিভাজনে
$\alpha=\beta$। অতএব থামার আগের অংশে
$g$ !!{injective}।

এবার দেখাই, কোনো ক্ষুদ্রতম ক্রমসংখ্যা $\alpha$-তে
সত্যিই থামতে হয়, অর্থাৎ
$A=g[\alpha]$। নইলে $g$ সর্বত্র
!!{injective} হতো; আগের
যুক্তিতে প্রত্যেক ক্রমসংখ্যা $\alpha$-র জন্য
% BN-SRC-469: an injection is into A, not Pow(A) minus the empty set;
% Hartogs contradicts non-strict embeddability for every ordinal.
$\cardle{\alpha}{A}$ হতো, যা
\olref[hartogs]{HartogsLemma}-এর বিরোধ।
থামার আগের অংশের বিস্তার ঠিক $A$।

সব মিলিয়ে, $g$-র সেই প্রারম্ভিক অংশটি
কোনো ক্রমসংখ্যা থেকে $A$-তে একটি
!!a{bijection}। $g$-র এই অংশ দিয়ে $A$
সুক্রমবিন্যস্ত করা যায়।
\end{proof}

অতএব সুক্রমবিন্যাস ও চয়ন একসঙ্গে প্রতিষ্ঠিত হয়,
অথবা একসঙ্গে ব্যর্থ হয়। কিন্তু প্রশ্ন রয়ে যায়:
এরা প্রতিষ্ঠিত হয়, না ব্যর্থ?

\end{document}
